\VerdierPhysicalPageAnchor{0074}%
\VerdierSemanticAnchor{ch1-2-4-1-pi-opposite}
the functor opposite to the functor defined in (2.2.3) and (2.2.4).
Similarly, let us denote by:

\VerdierSemanticAnchor{ch1-2-4-1-eq4}
\begin{equation}
  \int_{\varphi,C^\circ}^{\Pi}\colon
  \operatorname{comp}^{n}(C^\circ)
  \longrightarrow
  \operatorname{comp}^{m}(C^\circ),
  \tag{2.4.1.4}
\end{equation}

\VerdierSemanticAnchor{ch1-2-4-1-pi-definition}
the functor defined in (2.2.3) and (2.2.4). One verifies immediately that
the diagram:

\VerdierSemanticAnchor{ch1-2-4-1-eq5}
\begin{equation}
  \resizebox{0.82\textwidth}{!}{%
    $\begin{tikzcd}[ampersand replacement=\&,column sep=6.5em,row sep=5.2em]
      \operatorname{comp}^{n}(C)^\circ
        \arrow[r,"\iota_n"]
        \arrow[d,"{\left(\int_{\varphi,C}^{\Sigma}\right)^\circ}"']
        \& \operatorname{comp}^{n}(C^\circ)
          \arrow[d,"{\int_{\varphi,C^\circ}^{\Pi}}"] \\
      \operatorname{comp}^{m}(C)^\circ
        \arrow[r,"\iota_m"]
        \& \operatorname{comp}^{m}(C^\circ)
    \end{tikzcd}$%
  }
  \tag{2.4.1.5}
\end{equation}

\VerdierSemanticAnchor{ch1-2-4-1-eq5-commutative}
is commutative. In particular, when $\varphi$ is the unique map
$[n]\longrightarrow[1]$, there is a commutative diagram:

\VerdierSemanticAnchor{ch1-2-4-1-eq6}
\begin{equation}
  \resizebox{0.82\textwidth}{!}{%
    $\begin{tikzcd}[ampersand replacement=\&,column sep=6.5em,row sep=5.2em]
      \operatorname{comp}^{n}(C)^\circ
        \arrow[r,"\iota_n"]
        \arrow[d,"{\left(\int_{n,C}^{\Sigma}\right)^\circ}"']
        \& \operatorname{comp}^{n}(C^\circ)
          \arrow[d,"{\int_{n,C^\circ}^{\Pi}}"] \\
      \operatorname{comp}(C)^\circ
        \arrow[r,"\iota_1"]
        \& \operatorname{comp}(C^\circ)
    \end{tikzcd}$%
  }
  \tag{2.4.1.6}
\end{equation}

\VerdierSemanticAnchor{ch1-2-4-1-twisted-isomorphism}
and one verifies immediately that the pair of isomorphisms
$(\iota_1,\iota_n)$ is an isomorphism from the twisted
$\mathbb Z[n]$-functor $\left(\int_{n,C}^{\Sigma}\right)^\circ$ onto the
twisted $\mathbb Z[n]$-functor $\int_{n,C^\circ}^{\Pi}$ (the graded
structure on $\int_{n,C^\circ}^{\Pi}$ is described in (2.3.4); the graded
structure on $\left(\int_{n,C}^{\Sigma}\right)^\circ$ is described in
(2.3.4) and (1.4.8)).

\VerdierSemanticAnchor{ch1-2-4-1-identification}
Hereafter, by means of the isomorphism $\iota_n$, we identify the
categories $\operatorname{comp}^{n}(C)^\circ$ and
$\operatorname{comp}^{n}(C^\circ)$, and the functors
$\left(\int_{\varphi,C}^{\Sigma}\right)^\circ$ and
$\int_{\varphi,C^\circ}^{\Pi}$.

\VerdierSemanticAnchor{ch1-2-4-2}
\noindent\textbf{2.4.2.}\quad
Let $C_i$, $i\in[n]$, be a family of additive categories, and let
$\prod_{i\in[n]}C_i$ be their product. We denote the objects of
$\prod_{i\in[n]}C_i$ by $(X_1,X_2,\ldots,X_n)$ and use the notation
$(X_1,X_2,\ldots,f_i,\ldots,X_n)$ to denote the morphism:

\VerdierSemanticAnchor{ch1-2-4-2-component-morphism}
\begin{equation*}
  (\operatorname{id}_{X_1},\operatorname{id}_{X_2},\ldots,
  f_i,\ldots,\operatorname{id}_{X_n}).
\end{equation*}
